La reacció de fusió més simple és la fusió d'un protó amb un neutró. El resultat d'aquesta fusió és la formació d'un determinat isòtop de l'hidrogen i d'un fotó.

\begin{apartat}{1}
Quin isòtop es forma? Escriviu l'equació nuclear que correspon a aquest procés.
\end{apartat}

\begin{apartat}{1}
Determineu l'energia del fotó en joules (J) i en electró-volts (eV). Calculeu la freqüència d'aquest fotó.
\end{apartat}

\dades{%
  Isòtops de l'hidrogen: ${}^{1}_{1}\mathrm{H}$, ${}^{2}_{1}\mathrm{H}$, ${}^{3}_{1}\mathrm{H}$.\\
  Velocitat de la llum, $c = 3,00\times 10^{8}\ \mathrm{m\,s^{-1}}$.\\
  Constant de Planck, $h = 6,63\times 10^{-34}\ \mathrm{J\,s}$.\\
  $1\ \mathrm{eV} = 1,602\times 10^{-19}\ \mathrm{J}$.\\
  Masses en repòs:\\[3pt]
  {\small\begin{tabular}{|c|c|c|c|}
  \hline
  ${}^{1}_{1}\mathrm{H}$ (protó) & ${}^{1}_{0}\mathrm{n}$ (neutró) & ${}^{2}_{1}\mathrm{H}$ (deuteri) & ${}^{3}_{1}\mathrm{H}$ (triti)\\
  \hline
  $1,672\,62\times 10^{-27}\ \mathrm{kg}$ & $1,674\,92\times 10^{-27}\ \mathrm{kg}$ & $3,343\,58\times 10^{-27}\ \mathrm{kg}$ & $5,007\,36\times 10^{-27}\ \mathrm{kg}$\\
  \hline
  \end{tabular}}}
